\subsection{Scopo del documento}
Il presente documento costituisce il \textbf{Manuale Cliente} del sistema \textit{SmartOrder}, sviluppato dal gruppo NightPRO nell'ambito del corso di Ingegneria del Software (A.A.\ 2025/2026) dell'Universit\`a degli Studi di Padova.

Lo scopo del manuale \`e fornire agli utenti finali CLIENTI tutte le informazioni necessarie per utilizzare in modo efficace la piattaforma, descrivendone le funzionalit\`a, l'interfaccia e le modalit\`a di interazione.

Il documento \`e redatto in modo incrementale: le sezioni relative alle istruzioni d'uso sono state completate a seguito dello sviluppo dell'MVP.

\subsection{Scopo del prodotto}
\textit{SmartOrder} \`e una piattaforma B2B\textsubscript{\textit{g}} per la gestione automatizzata degli ordini commerciali. Il sistema consente ai clienti di formulare ordini in \textbf{linguaggio naturale} --- tramite testo e audio --- sfruttando un pipeline di intelligenza artificiale basato su LLM\textsubscript{\textit{g}} per:

\begin{itemize}[noitemsep]
    \item interpretare l'intento dell'utente e dedurre quantit\`a e prodotti;
    \item mappare gli articoli richiesti sul catalogo aziendale tramite \textbf{Vector Search} (embedding);
    \item generare un output strutturato in formato JSON\textsubscript{\textit{g}} pronto per l'integrazione con il sistema ERP\textsubscript{\textit{g}} aziendale.
\end{itemize}

La piattaforma mette inoltre a disposizione un \textbf{carrello virtuale} per la validazione visiva dell'ordine prima dell'invio, un \textbf{sistema di ticketing} per l'assistenza umana in tempo reale e un meccanismo di \textbf{logging e feedback} per il miglioramento continuo del modello (Human-in-the-Loop).

\subsection{Destinatari del documento}
Il manuale si rivolge alla seguente categoria di utenti:

\begin{description}[style=nextline]
    \item[\textbf{Cliente}] Acquirente B2B (ad esempio ristoratore, negoziante) che utilizza l'interfaccia mobile-first per dettare ordini, gestire il carrello e consultare lo storico.
\end{description}

\subsection{Glossario}
Per evitare ambiguit\`a, i termini tecnici e di dominio sono definiti nel \textit{Glossario v2.0}.
I termini presenti nel glossario sono contrassegnati con il simbolo \textsubscript{\textit{g}} alla prima occorrenza in ogni sezione.

\subsection{Riferimenti}

\subsubsection{Riferimenti normativi}
\begin{itemize}
    \item \href{https://swenightpro.github.io/Documentazione/docs/PB/Documentazione\%20Interna/Norme_Di_Progetto.pdf}{Norme di Progetto v4.0 (Ultima modifica: 2026-04-09)}
    \item \href{https://swenightpro.github.io/Documentazione/docs/PB/Documentazione\%20Esterna/Analisi_Dei_Requisiti.pdf}{Analisi dei Requisiti v3.0 (Ultima modifica: 2026-04-09)}
    \item \href{https://swenightpro.github.io/Documentazione/docs/PB/Documentazione\%20Esterna/Specifica_Tecnica.pdf}{Specifica Tecnica v1.0 (Ultima modifica: 2026-04-09)}
\end{itemize}

\subsubsection{Riferimenti informativi}
\begin{itemize}
    \item \href{https://swenightpro.github.io/Documentazione/docs/PB/Documentazione\%20Esterna/Glossario.pdf}{Glossario v3.0 (Ultima modifica: 2026-04-09)}
    \item \href{https://nextjs.org/docs}{Next.js Documentation (Ultima modifica: 2026-04-09)}
    \item \href{https://fastapi.tiangolo.com}{FastAPI Documentation (Ultima modifica: 2026-04-09)}
\end{itemize}